%% sections/conclusion.tex -- Phase 7.1. Register: CLAUDE.md S9.2.

\section{Conclusion}
\label{sec:conclusion}

An instructor adopting an LLM grader faces several choices, not one, and prior studies addressed them
only in fragments, with some left unstudied. We aggregated those fragments, verified them for computer-science
grading under one controlled protocol, and condensed the result into a decision guide. Three results carry the weight. Reasoning is a
conditional trade, not an upgrade: it bought agreement in a minority of model-dataset cells, always
at a large, domain-conditional token premium and almost always at a consistency cost. Its one
reliably explained success was repairing a grader whose baseline had collapsed to zero on code.
Evaluation guidance is the cheapest reliable lever for the model we tested it on. And the configuration that ranked best on the
public data also led in deployment, so it makes a strong default, even though the verified samples were
too small to confirm the order below the top. Together these define the framework's stance: the right
configuration depends on the task, not on a single best grader, and by keying every prior to observable
properties rather than to model identity, we built the guidance to outlast the models we tested.

Future work follows from the gaps this study declares and from what the literature has yet to settle.
Two gaps now arise in our design. Few-shot guidance, adding a few already-graded answers to the prompt
as examples, was held out; the strong effect of the studied guidance makes it worth testing, but only with a
principled rule for choosing the examples. And a
fine-tuning arm remains open: fine-tuning is already competitive for grading~\cite{cong_memorization_2026},
but it has not been placed in this framework, comparing a fine-tuned grader against its base model on
agreement, consistency, and cost together. Three more follow from the
literature. Code evidence rests on two Java-centric datasets because human-rubric-graded code corpora are
scarce~\cite{havare_ai-based_2025,pathak_rubric_2025}, so widening it to other languages and paradigms,
and contributing such data, is the clearest way to firm up the code conclusions. Transfer was shown in a
single Portuguese deployment, so testing the priors across more institutions and student populations would
establish how far they carry. And the low-end score collapse we found on code, where a grader floors much
of the scale at zero, is not described in the grading literature and deserves study in its own right, both
to detect it and to prevent it without paying for reasoning.
